\documentclass{article}
\usepackage{tkz-tab}
\usepackage{amsmath}
\usepackage[french]{babel}
\usepackage[T1]{fontenc}

\title{Mathématiques lycée}
\author{JIVA LILA Rehan}
\date{2022\textemdash 2024}

\begin{document}

\maketitle
    Ce livret est un récapitulatif de mon année de mathématiques de première et terminale, au Lycée International de Ferney-Voltaire, sur le site principal, durant les années 2022\textemdash2024, dans la classe 1G10 et de TG7.

\tableofcontents

\pagebreak



\section{Fonctions de références}


\subsection{Propriétés de fonction}

\subsubsection{Croissance \& décroissance}
    \underline{Pour $a<b$:}\\
    Si $f(a) < f(b)$: la fonction est \emph{croissante}\\
    Si $f(a) > f(b)$: la fonction est \emph{décroissante}\\
    Si $f(a) = f(b)$: la fonction est \emph{monotone}

\subsubsection{Parité \& imparité}
    \underline{$\forall x \in \mathbf{R}$ :}\\
    Si $f(-x) = f(x)$: la fonction est \emph{paire}\\
    Si $f(-x) = -f(x)$: la fonction est \emph{impaire}


\subsection{Fonctions affines}
    Une fonction est une fonction de style $f(x) = ax+b$ avec $a$ la dérivée de la fonction, et $b$ son ordonnée à l'origine.


\subsection{Fonctions carrées}
    La fonction carrée est définie par $f(x) = x^2$. Elle est paire sur $\mathbf{R}$, décroissante sur $\mathbf{R^-}$ et croissante sur sur $\mathbf{R^+}$.


\subsection{Fonction cube}
    La fonction cube est définie par $f(x) = x^3$. Elle est impaire et croissante sur $\mathbf{R}$.


\subsection{Fonction inverse}
    La fonction inverse est définie par $f(x) = \frac{1}{x}$. Elle est impaire et décroissante sur $\mathbf{R^*}$.


\subsection{Fonction racine carrée}
    La fonction carrée est définie sur $\mathbf{R^+}$ par $f(x) = \sqrt{x} = x^{-2}$. Elle est croissante.


\subsection{La fonction valeur absolue}
    La fonction valeur absolue est définie par $f(x) = |x| $. Elle est paire sur sur $\mathbf{R}$, décroissante sur $\mathbf{R^-}$ et croissante sur sur $\mathbf{R^+}$.\\
    Son comportement suit la règle suivante : $|x| = |-x| = x$.



\section{Règles sur les inégalités}
    \begin{enumerate}
        \item Les additions et soustraction \emph{ne changent pas le sens}.
        \item Les multiplication et division \emph{ne changent le sens QUE SI leur produit ou dividende est négatif}.
        \item L'application d'une fonction \emph{ne change le sens QUE SI sa fonction est impaire}.
    \end{enumerate}



\section{Fonctions polynômes du second degré}


\subsection{Formes}

\subsubsection{Forme développée (ou forme développée réduite)}
$f(x) = ax^2+bx+c$ avec $a \neq 0$

\subsubsection{Forme factorisée}
$f(x) = a(x - x_1)(x - x_2)$ avec $x_1$ et $x_2$ les racines de $f$\\
On notera que $x_1+x_2=-b/a$ et que $x_1x_2=c/a$

\subsubsection{Forme canonique}
$f(x) = a(x-\alpha)+\beta$ avec $\alpha = \frac{-b}{2a}$ et $\beta = f(\alpha)$


\subsection{Variations}
    Pour $a > 0$: décroissante avant $\alpha$, avec un minimum en $\beta$, puis croissante.\\
    Pour $a < 0$: croissante avant $\alpha$, avec un maximum en $\beta$, puis décroissante.


\subsection{Résolution}
    \begin{enumerate}
        \item On calcule le discriminant tel que $\Delta = b^2 - 4ac$.
        \item On calcule les racines telles que $x_0 = \frac{-b \pm \sqrt{\Delta}}{2a}$.
    \end{enumerate}


\subsection{Signe de la fonction}
    \begin{enumerate}
        \item Si $\Delta > 0$: Du signe de $a$ à l'extérieur des racines, de l'autre entre elles. S'annule sur ces dernières
        \item Si $\Delta = 0$: Du signe de $a$ sur $\mathbf{R}$ S'annule sur les racines.
        \item Si $\Delta < 0$ Du signe de $a$ sur $\mathbf{R}$.
    \end{enumerate}



\section{Dérivations}


\subsection{Dérivation locale}

\subsubsection{Taux de variation}
    \begin{enumerate}
        \item Entre $a$ et $b$ : $\frac{f(b) - f(a)}{b - a}$.
        \item Sur $a$: $\frac{f(a + h) - f(a)}{h}$.
    \end{enumerate}

\subsubsection{Nombre dérivé (dérivée)}
    \textit{La dérivée de $f(x)$, soit la variation de $f$ en $x$, s'écrit $f'(x)$.}
    $$f'(x) = \lim \limits_{h \Rightarrow 0} \frac{f(a + h) - f(a)}{h}$$

\subsubsection{Tangente à une courbe}
    La tangente à la courbe de $f$ en $a$ est définie par: $$\mathcal{T}(x) = f'(a)(x-a) + f(a)$$


\subsection{Dérivation globale}
    Une dérivation globale reprend le concept de dérivation locale, mais fourni une formule applicable sur toute la fonction. Elle est utilisée pour étudier les variations de la fonction.

\subsubsection{Opérations sur les dérivées}
    \begin{enumerate}
        \item[] $$(u+v)' = u' +v'$$
        \item[] $$\forall k \in \mathbf{R} : (ku)' = ku'$$
        \item[] $$(uv)' = vu' + uv'$$
        \item[] $$\frac{1}{v}' = \frac{-v'}{v^2}$$
        \item[] $$\frac{u}{v}' = \frac{u'v - v'u}{v^2}$$
        \item[] À savoir: $\forall f = g(h) : f' = h'g(h)$ avec $f$, $g$, et $h$ des fonctions
    \end{enumerate}

\subsubsection{Dérivées de fonctions usuelles}
    \begin{enumerate}
        \item[] $$\forall k \in \mathbf{R}, \forall f(x) = k : f'(x) = O$$
        \item[] $$\forall k \in \mathbf{R}, \forall f(x) = kx : f'(x) = k$$
        \item[] $$\forall n, \forall f(x) = x^n : f'(x) = nx^{n-1}$$
        \item[] $$\forall x \in \mathbf{R^*}, \forall f(x) =\frac{1}{x} : f'(x) = \frac{-1}{x^2}$$
        \item[] $$\forall x \in \mathbf{R^{+*}}, \forall f(x) =\sqrt{x} : f'(x) = \frac{1}{2\sqrt{x}}$$
        \item[] $$\forall f(x) = e^{g} : f'(x) = g'e^{g}$$
        \item[] $$\forall x \in \mathbf{R^{-*}}, \forall f(x) = |x| : f'(x) = -1$$
        \item[] $$\forall x \in \mathbf{R^{+*}}, \forall f(x) = |x| : f'(x) = 1$$
    \end{enumerate}

\subsubsection{Variations de la fonction}
    $\forall f$, si $f'$ est négatif, $f$ est décroissante. Inversement, si $f'$ est positive, $f$ est croissante.



\section{Suites numériques}


\subsection{Définition}
    Une suite numérique $u$ est une suite de valeurs pourvues d'un indice, $n$, qui parcours l'ensemble $\mathbf{N}$. L'élément à l'indice $n$ de la suite $u$ est noté $u_n$.


\subsection{Modes de génération}

\subsubsection{Liste}
    Une simple liste de nombres. (ex: les décimales de $\pi$: 1; 4; 1; 5\dots)

\subsubsection{Formule explicite}
    Définir $u_n$ en fonction de $n$: $u_n = f(n)$

\subsubsection{Récurrence}
    Définir $u_n$ en fonction de un ou plusieurs $u_{n-k}$ (ex: suite de Fibonacci: $u_n = u_{n-2} + u_{n1}$)


\subsection{Quelques styles de suites}

\subsubsection{Suite arithmétique}
    Une suite arithmétique est définie par $u_n = u_{n-1} + r$ ou $u_n = u_0 + nr$ avec $r$ sa raison.\\
    La somme de ses termes d'indices de $p$ à $n$ vaut $(n+1)\times \frac{u_p+u_n}{2}$.\\
    Pour prouver qu'une suite est arithmétique, il suffit que $r$ dans $u_n - u_{n-1} = r$ soit indépendant de $n$.

\subsubsection{Suite géométrique}
    Une suite arithmétique est définie par $u_n = u_{n-1} \times q$ ou $u_n = u_0 \times q^n$ avec $q$ sa raison.\\
    La somme de ses termes d'indices de $p$ à $n$ vaut $u_p\times \frac{1-q^{n-p+1}}{1-q}$.\\
    Pour prouver qu'une suite est géométrique, il suffit que $q$ dans $\frac{u_n}{u_{n-1}} = q$ soit indépendant de $n$.


\subsection{Croissance et décroissance}
    Pour prouver la croissance ou la décroissance d'une suite, on applique \\$u_{n-+1} - u_n$, qui si négatif prouve une décroissance, et inversement.


\subsection{Limite de suites}

\subsubsection{Définition}
    \begin{enumerate}
        \item $v_n$ admet une limite $+\infty$ si $\forall A \in \mathbf{R}$, il existe un rang $n_0$ tel que $\forall n \ge n_0$, $u_n > A$
        \item $v_n$ admet une limite $-\infty$ si $\forall A \in \mathbf{R}$, il existe un rang $n_0$ tel que $\forall n \ge n_0$, $u_n < A$
        \item $v_n$ admet une limite $l$ si $\forall E \in \mathbf{N^*}$, il existe un rang $n_0$ tel que $\forall n \ge n_0$ $l-E<u_n<l+E$
    \end{enumerate}

\subsubsection{Opérations}
    Lors du calcul d'une limite, les formes indéterminées sont:
    \\$$\infty - \infty \qquad 0 \times \infty \qquad \frac{0}{0} \qquad \frac{\infty}{\infty}$$
    Si on en rencontre une, l'usage est de manipuler l'expression de sorte à ne plus avoir ce problème, ou bien d'appliquer la règle de L'Hôpital qui dit que $$\lim \limits_{x \Rightarrow e} \frac{f(x)}{g(x)} = \lim \limits_{x \Rightarrow e} \frac{f'(x)}{g'(x)}$$ ssi $f'(x)$ et $g'(x)$ existent et $g'(x) \ne 0$.

\subsubsection{Comparaison et encadrement}
    On suppose deux suites $u_n$ et $v_n$ telles que après un certain rang, $u_n>v_n$.
    $$\lim \limits_{n \Rightarrow \infty} v_n = + \infty \Rightarrow \lim \limits_{n \Rightarrow \infty} u_n = + \infty$$
    $$\lim \limits_{n \Rightarrow \infty} u_n = - \infty \Rightarrow \lim \limits_{n \Rightarrow \infty} v_n = - \infty$$
    On suppose $u_n$, $v_n$ et $w_n$, telles que après un certain rang, $u_n \le v_n \le w_n$
    $$\lim \limits_{n \Rightarrow \infty} u_n = \lim \limits_{n \Rightarrow \infty} w_n = l \Rightarrow \lim \limits_{n \Rightarrow \infty} v_n = l$$



\section{Trigonométrie}


\subsection{Cercle trigonométrique}
    Le cercle trigonométrique est défini par un cercle de rayon $1 UA$, donc de rayon $2\pi$.\\
    On admettra que son centre $O$ est aux coordonnées $(0;0)$.


\subsection{Radian}
    Le radian est l'unité de mesure d'angle en trigonométrie. Note $rad$, on retiendra que $1^\circ= \frac{\pi}{180} rad$


\subsection{Points sur l'arc}
Le point correspondant a l'angle $x$ radians, aura pour coordonnées $(\cos{x};\sin{x})$. \\De plus, l'angle $x$ est le même que l'angle $x +2k\pi$, $\forall k \in \mathbf{N}$.


\subsection{Fonctions cosinus et sinus}
    \begin{enumerate}
        \item $\cos x$ est une fonction paire, variant entre $-1$ et $1$, et récurrente telle que $\cos{x + 2k\pi} = \cos x$, $\forall k \in \mathbf{N}$.
        \item $\sin x$ est une fonction impaire, variant entre $-1$ et $1$, et récurrente telle que $\sin{x + 2k\pi} = \sin x$, $\forall k \in \mathbf{N}$.
    \end{enumerate}



\section{Probabilités}


\subsection{Vocabulaire}
    Soit deux évènements $\mathcal{A}$ et $\mathcal{B}$ :\\
    On note $\mathcal{A} \cup \mathcal{B}$ l'union des deux évènements, soit celui ou au moins l'un des deux se produit.\\
    On note $\mathcal{A} \cap \mathcal{B}$ l'intersection des deux évènements, soit celui ou les deux se produisent.\\
    On note $\mathcal{_AB}$ l'évènement ou $\mathcal{B}$ se produit, sachant que $\mathcal{A}$ s'est déjà produit (utilise typiquement dans les arbres pondérés).\\


\subsection{Formules}
    $$\mathcal{P_A} \cup \mathcal{P_B} = \mathcal{P_A} +\mathcal{P_B} - \mathcal{P_A} \cap \mathcal{P_B}$$
    $$\mathcal{P_A} \cap \mathcal{P_B} = \mathcal{P_A} +\mathcal{P_B} - \mathcal{P_A} \cup \mathcal{P_B}$$
    $$\mathcal{P_AB} = \frac{\mathcal{PA} \cap \mathcal{PB}}{\mathcal{PA}}$$
    $$\mathcal{P_AB} = \frac{\mathcal{P_BA} \times \mathcal{PB}}{\mathcal{PA}}$$


\subsection{Indépendance}
    $\mathcal{A}$ et $\mathcal{B}$ sont indépendants si et seulement si $\mathcal{PA} \cap \mathcal{PB} = \mathcal{PA} \times \mathcal{PB}$. \\
    Des lors, $\mathcal{P_BA} = \mathcal{PA}$ et $\mathcal{P_AB} = \mathcal{PB}$.



\section{Variables aléatoires}


\subsection{Définition}
    Une variable aléatoires est une variable pouvant prendre différentes valeurs, suivant une certaine probabilité. \\
    On associe alors à la variable aléatoire $X$ les valeurs $x_1, x_2, ..., x_n$ et leurs probabilités respectives $p_i$ définies par $X == x_i$.


\subsection{Paramètres}
    Une variable aléatoire a dès lors des paramètres, définis comme suis:
    \begin{enumerate}
        \item[] L'espérance,$E(X)$, qui représente la moyenne pondérée des valeurs possibles de cette variable, et qui est définie par: $$\sum^n_{i=1}p_ix_i$$
        \item[] Sa variance, $V(X)$: $$\sum^n_{i=1}p_i(x_i-E(X))^2$$
        \item[] Son écart-type, $\sigma(X)$: $$\sqrt{V(X)}$$
    \end{enumerate}



\section{Géométrie repérée}


\subsection{Vecteurs}
    Propriétés:
    \begin{enumerate}
        \item On note la norme $AB$ du vecteur $\Vec{AB}$ : $||\Vec{AB}||$.
        \item Le vecteur opposé a $\Vec{AB}$ est $-\Vec{AB} = \Vec{BA}$.
        \item $\Vec{AB}+\Vec{BC}=\Vec{AC}$
        \item Si $\Vec{v} = k\Vec{u}$ avec $k$ un réel, alors $\Vec{v}$ et      $\Vec{u}$ sont colinéaires.
        \item Dans un repère orthonormé, $||\Vec{AB}|| = \sqrt{(x_B-x_A)^2+ (y_B-y_A)^2+...}$.
    \end{enumerate}

\subsubsection{Produit scalaire}
    Le produit scalaire, ou travail, d'un vecteur sur un autre, se note $\mathcal{W}$.\\
    On notera que $\mathcal{W}_uv = \mathcal{W}_vu$.\\
    Voici les méthodes de calcul de $\Vec{AB} \cdot \Vec{AC} = \mathcal{W}_{AB}AC$.\\
    \\
    Méthodes de calcul:
    \begin{enumerate}
        \item Par le projeté orthogonal \\Soit $H$ le projeté orthogonal de $C$ sur $(AB)$: $\Vec{AB} \cdot \Vec{AC} = \Vec{AB} \cdot \Vec{AH} = AB \times AH$.
        \item Par les coordonnées \\$\Vec{AB} \cdot \Vec{AC} = x_{AB}x_{AC}+y_{AB}y_{AC}$
        \item Par le cosinus \\$\Vec{AB} \cdot \Vec{AC} = AB\times AH\times \cos \alpha$ avec $\alpha$ l'angle forme par les vecteurs.
        \item Par les normes \\$\Vec{AB} \cdot \Vec{AC} = \frac{1}{2}(AB^2 + AC^2 - BC^2)$
    \end{enumerate}
    On notera que $\Vec{AB}$ et $\Vec{AC}$ sont orthogonaux si leur produit scalaire est nul.


\subsection{Plan}
    Un plan est défini dans l'espace par deux vecteurs non colinéaires et un point.\\
    Dès lors pour un plan $\mathcal{P}$ défini par le point $A$ et les vecteurs $\Vec{u}$ et $\Vec{v}$, $M \in \mathcal{P}$ ssi $\Vec{AM} = a\Vec{u}+b\Vec{v}$ avec $a,b\in\mathbf{R}$

\subsection{Base de l'espace}
    Trois vecteurs non coplanaires forment une base de l'espace\\
    Dès lors, $\forall \Vec{u} \in$ la base de l'espace $(\Vec{i};\Vec{j};\Vec{k})$, $\exists x,y,z \in \mathbf{R}$ tels que $x\Vec{u}+y\Vec{j}+z\Vec{k}=\Vec{u}$. On dira que $\Vec{i} = \begin{pmatrix}x\\y\\z\end{pmatrix}$ dans $(\Vec{i};\Vec{j};\Vec{k})$\\\\
    Propriétés:
    \begin{enumerate}
        \item $\Vec{AB} = (x_b - x_a ; y_b - y_a ; ...)$
        \item $I$ est milieu de $AB$ ssi $I=(\frac{x_a+x_b}{2};\frac{y_a+y_b}{2};...)$
    \end{enumerate}


\subsection{Formes géométriques}

\subsubsection{Droite}
    Une équation de droite peut s'écrire de deux manières: la forme réduite $y=mx+p$, et la forme cartésienne $ay+bx+c=0$.\\
    On distingue alors le vecteur directeur et le vecteur normal (orthogonal au premier), de formules respectives $\begin{pmatrix}-b\\a\end{pmatrix}$ et $\begin{pmatrix}a\\b\end{pmatrix}$. On peut alors exploiter ces vecteurs pour déterminer des droites parallèles ou perpendiculaires à la première.
    \\\\
    Ajout en terminale:\\
    Soit $(\Vec{i};\Vec{j};\Vec{k})$ une base de l'espace et $D$ une droite $(x_a;y_a;z_a;\Vec{\begin{pmatrix}a\\b\\c\end{pmatrix}})$,\\ $M(x_m;y_m;z_m) \in D$ ssi  \[
                                \left \{
                                \begin{array}{c @{=} c}
                                    x_m & x_a + ak \\
                                    y_m & y_a + bk \\
                                    z_m & z_a + ck
                                \end{array}
                                \right.
                                \]
    avec $k \in \mathbf{R}$

\subsubsection{Cercle}
    Ici la forme cartésienne est $(x-x_A)^2+(y-y_A)^2 = r^2$ avec $A$ le centre du cercle et $r$ son rayon.\\
    Dès lors, l'ensemble des les points $(x;y)$ qui vérifient $(x-a)^2+(y-b)^2 = S$ avec $S\in\mathbf{R}$, est un cercle de centre $(a;b)$ et de rayon $\sqrt{S}$.

\subsubsection{Parabole}
    On retrouve la forme $y=ax^2+bx+c$ où on remarque l'axe de symétrie $x = -b/2a$ et le sommet $(-b/2a,f(-b/2a))$.



\section{Boite à outils}


\subsection{Prouver une proposition}

\subsubsection{Par déduction}
    On part d'une hypothèse et par un raisonnement logique, on arrive à une conclusion.

\subsubsection{Par contre exemple}
    On trouve un cas qui contredit la proposition.

\subsubsection{Par l'absurde}
    On suppose P et arrive à une contradiction.

\subsubsection{Par contraposée}
    On prouve la contraposée pour prouvée la proposition

\subsubsection{Par disjonction de cas}
    Lorsque la proposition dépend d'un variable, on sépare le raisonnement suivant les valeurs de la dite variable.

\subsubsection{Par récurrence}
    On prouve qu'une proposition est vrai à partir d'un certain rang, en prouvant qu'elle est initialisée pour ce rang puis qu'elle est héréditaire.




\end{document}
